\documentclass[11pt,a4paper]{article}

\usepackage[a4paper,left=2.0cm,right=2.0cm,top=1.55cm,bottom=1.55cm]{geometry}
\usepackage[T1]{fontenc}
\usepackage{lmodern}
\usepackage{graphicx}
\usepackage{xcolor}
\usepackage{array}
\usepackage{tabularx}
\usepackage{ragged2e}
\usepackage{enumitem}
\usepackage{fancyhdr}
\usepackage{tikz}
\usepackage{setspace}


% -------------------------------------------------
% COLOURS
% -------------------------------------------------
\definecolor{TitleBlue}{RGB}{61,103,154}

% -------------------------------------------------
% GENERAL SETTINGS
% -------------------------------------------------
\setlength{\parindent}{0pt}
\setlength{\parskip}{0pt}
\renewcommand{\arraystretch}{1.0}

\newcommand{\blueheading}[1]{%
    {\sffamily\color{TitleBlue}\fontsize{15}{18}\selectfont #1}%
}

\newcommand{\bluebigheading}[1]{%
    {\sffamily\bfseries\color{TitleBlue}\fontsize{16}{19}\selectfont #1}%
}

\newcommand{\instruction}[1]{%
    {\itshape\fontsize{10.5}{12.5}\selectfont #1}%
}

\newcommand{\answerbox}[2][3.0cm]{%
    \vspace{2pt}
    \noindent
    \fbox{%
        \begin{minipage}[t][#1][t]{\dimexpr\textwidth-2\fboxsep-2\fboxrule\relax}
            \vspace{1mm}
            \itshape #2
        \end{minipage}
    }
}

% -------------------------------------------------
% HEADER / FOOTER
% -------------------------------------------------
\pagestyle{fancy}
\fancyhf{}
\renewcommand{\headrulewidth}{0pt}
\renewcommand{\footrulewidth}{0pt}

\fancyhead[L]{%
    \sffamily\bfseries\itshape\fontsize{10.5}{12}\selectfont
    Data Structures with C(TCS302)

    Object Oriented Programming with c++(TCS307)
}
\fancyhead[R]{%
    \sffamily\bfseries\fontsize{10.5}{12}\selectfont
    PROJECT PROPOSAL
}
\fancyfoot[R]{%
    \itshape\fontsize{10}{12}\selectfont
    Page \thepage
}
\setlength{\headheight}{14pt}
\setlength{\headsep}{23pt}
\setlength{\footskip}{24pt}

% -------------------------------------------------
% PAGE 1
% -------------------------------------------------
\begin{document}

\vspace*{4mm}

\bluebigheading{PROJECT AND TEAM INFORMATION}

\vspace{13mm}

\blueheading{Project Title}

\vspace{1mm}


\answerbox[1.0cm]{GridRescue: Interactive Grid Failure \& Load Redistribution Simulator}

\vspace{13mm}

\blueheading{Student / Team Information}

\vspace{2mm}

\noindent
\begin{tabularx}{\textwidth}{|>{\RaggedRight\arraybackslash}p{0.69\textwidth}|>{\RaggedRight\arraybackslash}X|}
\hline
\begin{minipage}[t]{\linewidth}
\vspace{1.5mm}
\textit{Team Name:}\\[-1mm]
\textit{Team \#}
\vspace{1.5mm}
\end{minipage}
&
\begin{minipage}[t]{\linewidth}
\vspace{1.5mm}
\textit{GridRescuers}\\
\textit{ DSCPP-III-2026-T091}
\vspace{1.5mm}
\end{minipage}
\\
\hline
\begin{minipage}[t][3.25cm][t]{\linewidth}
\vspace{1.5mm}
\textbf{Team member 1 (Team Lead)}\\[-1mm]
\vspace{2.0cm}
\end{minipage}
&
\begin{minipage}[t][3.75cm][t]{\linewidth}
\vspace{1.5mm}
\textit{Nagarkoti, Tejasv -- \\ID: 2510010772}\\
\textit{tejasvsinghnagarkoti@gmail\\.com}

\vspace{2mm}
\includegraphics[width=3.0cm,height=1.5cm,keepaspectratio]{member1.jpeg}
\end{minipage}
\\
\hline
\begin{minipage}[t][3.25cm][t]{\linewidth}
\vspace{1.5mm}
\textbf{Team member 2}\\[-1mm]

\vspace{2.0cm}
\end{minipage}
&
\begin{minipage}[t][3.75cm][t]{\linewidth}
\vspace{1.5mm}
\textit{Vashisth,Tanishq -- \\ID: 2510010163}\\
\textit{tanishqvashisth10@gamil.\\com}

\vspace{2mm}
\includegraphics[width=3.0cm,height=1.5cm,keepaspectratio]{Report/member2.jpeg}
\end{minipage}
\\
\hline
\begin{minipage}[t][3.75cm][t]{\linewidth}
\vspace{1.5mm}
\textbf{Team member 3}\\[-1mm]
\vspace{2.0cm}
\end{minipage}
&
\begin{minipage}[t][3.25cm][t]{\linewidth}
\vspace{1.5mm}
\textit{Sahu,Tanishka -- \\ID:2510017468 }\\
\textit{Tanishka8874@gmail.com}

\vspace{2mm}
\includegraphics[width=3.0cm,height=1.5cm,keepaspectratio]{Report/member3.jpeg}
\end{minipage}
\\
\hline
\begin{minipage}[t][3.75cm][t]{\linewidth}
\vspace{1.5mm}
\textbf{Team member 4}\\[-1mm]
\vspace{2.0cm}
\end{minipage}
&
\begin{minipage}[t][3.25cm][t]{\linewidth}
\vspace{1.5mm}
\textit{Sandilya,Riddhiman -- \\ID:2510370168}\\
\textit{2007karmy@gmail.com}

\vspace{2mm}
\includegraphics[width=3.0cm,height=1.5cm,keepaspectratio]{Report/member4.jpeg}
\end{minipage}
\\
\hline
\end{tabularx}

\newpage

% -------------------------------------------------
% PAGE 2
% -------------------------------------------------

\bluebigheading{PROPOSAL DESCRIPTION }

\vspace{12mm}

\blueheading{Motivation }

\vspace{1mm}


\answerbox[6.5cm]{
\begin{itemize} 
\item Power grids often face unexpected breakdowns that cut off the electricity. When a fault happens, the first step is to figure out exactly what went wrong and which areas got disconnected.
\item After finding the issue, the system has to quickly look for another route to keep the power flowing and share whatever power is left among the users.
\item Through GridRescue, we are creating a basic simulation tool that catches these failures, finds a backup path, and manages the remaining load to keep the grid running.
\item We picked this project because it is a great way to use standard DSA and C++ OOPs concepts on a real-world problem instead of just solving basic code problems.
\end{itemize}
}

\vspace{8mm}
\vspace*{1.0cm}     
\blueheading{State of the Art / Current solution}

\vspace{1mm}


\answerbox[5.5cm]{
\begin{itemize}
\item Real-world power grids use large-scale systems like SCADA to monitor electricity flow. Whenever a fault happens, these systems automatically cut off the damaged part and find another way to send power.
\item But setting this up in real life is very expensive and hard. It needs thousands of sensors and 24x7 internet connectivity to process live data continuously.
\item In our GridRescue project, we are taking the basic idea behind these big systems and making a simple software version of it. In our app, a user can manually break a connection, and the program will use algorithms to find an alternate route.
\end{itemize}
}

\vspace{8mm}
\vspace*{1.0cm}     
\blueheading{Project Goals and Milestones }

\vspace{1mm}


\answerbox[5.5cm]{
The main goal of GridRescue is to build a small power grid simulation. We want to show what actually happens behind the scenes when a power line breaks. Users will see which areas face a blackout and how the system tries to restore power using other lines.

Our main milestones are:
\begin{itemize}
\item Represent the entire grid as a graph with power stations and consumers as nodes.
\item Figure out exactly which area got disconnected when a fault occurs.
\item Find a backup route to send power to those disconnected users.
\item Distribute whatever power is left properly among the users.
\end{itemize}
}

\vspace{12mm}
\newpage

\blueheading{Project Approach }

\vspace{1mm}



\answerbox[6cm]{
We will use a graph to design the power grid. Here, power stations and consumers will act as nodes, and the power lines will be stored using an Adjacency List. Whenever a user creates a fault, the program will run BFS and DFS to check which nodes (consumers) lost their connection.

After that, we will use Dijkstra's algorithm to find the shortest alternate path to restore the power. If there is a power shortage, we will use a Priority Queue to give electricity to important places like hospitals first. 

Everything will be coded in C++ using classes and vectors. We are also adding a simple GUI so anyone can easily click, break a wire, and see how the algorithm fixes it.
}



% -------------------------------------------------
% PAGE 3
% -------------------------------------------------
\vspace*{1.5cm}     
\blueheading{System Architecture (High Level Diagram)}

\vspace{3mm}




\answerbox[11.0cm]{
\begin{center}
\includegraphics[height=10.0cm]{sample1.jpeg}
\end{center}
}

\newpage

\vspace{12mm}

\blueheading{Project Outcome / Deliverables }

\vspace{1mm}



\answerbox[6.2cm]{
\begin{itemize}
\item By the end, we will submit a working GridRescue software with a simple graphical interface. It will show a basic map of power stations and consumers.
\item Users can create a failure in the system and see how the graph algorithms find the blackout areas in real-time.
\item The screen will update to show the new alternate paths and how much power each consumer is getting.
\item Along with the code, we will also provide a project report. This report will explain the working of our algorithms and how we used OOPs and DSA concepts to build this.
\end{itemize}
}

\vspace{12mm}

\blueheading{Assumptions}

\vspace{6mm}


\answerbox[3.5cm]{
To keep things simple, our grid will have a fixed size and fixed number of connections. The power capacity of substations and the demand of consumers will also be fixed. We will only test one failure at a time. The project will run on custom sample data, not real live data from a power company. Finally, during a power cut, certain places (like hospitals) will always get first preference.
}

\vspace{12mm}


\blueheading{References}

\vspace{1mm}

\instruction{}

\answerbox[5cm]{
1. Thomas H. Cormen et al., Introduction to Algorithms — for Graphs, BFS, DFS, Dijkstra and Priority Queue concepts.

2. E. Balagurusamy, Object Oriented Programming with C++ — for C++ classes, objects and basic OOP concepts.

3. C++ Documentation — for basic C++ concepts such as vector, priority\_queue and other standard library functions.

4. Power System Basics / Smart Grid Study Material — for understanding power grids, substations, consumers and power distribution.

5. Dear ImGui Documentation — for learning how to create the simple graphical interface for the project.
}

\end{document}
